% Combined online loop, checked against prefetch_policy.py and
% prefetch_training_benchmark.py. The algorithm describes one training epoch.
\begingroup
\color{blue}
\subsubsection{Combining Feature Caching and Reordering}
\label{subsec:caching-reorder-algorithm}

Algorithm~\ref{alg:caching-reorder} combines the two designs for one epoch.
The original sampler provides the batch order and the proxy for each batch.
The cache $C$ holds features before the first trainable layer, indexed by batch ID; $i$ is the next batch whose remaining training work must run.
The cache starts empty in each epoch, while the trainable parameters and loss history carry over.

\textsc{TrainRest} runs from the first trainable layer using the current parameters.
It keeps the original loss-based update decision, EMA update, and sampler feedback in Algorithm~\ref{alg:on-device_training}.
Its runtime estimate reserves time for a full backward update.
\textsc{FrozenForward} computes only the frozen layers before this point and returns features for each original batch.
Both runtime estimates are updated after the corresponding operation finishes.

\begin{algorithm}[t]
\color{blue}
\footnotesize
\SetAlFnt{\footnotesize\color{blue}}
\SetAlCapFnt{\footnotesize\color{blue}}
\SetAlCapNameFnt{\footnotesize\color{blue}}
\SetNlSty{textbf}{}{}
\SetAlgoNlRelativeSize{-1}
\SetAlgoNoEnd
\SetKw{Continue}{continue}
\SetKwInOut{Input}{Input}
\Input{Ordered batches $Q=(X_j,y_j,p_j)_{j=1}^{M}$; cache capacity $K$ in batches; lookahead $L$; allowed group sizes $S$; safety margin $\delta$; runtime estimates}
\BlankLine
$C\gets\emptyset$; $i\gets 1$\;
\While{$i\le M$}{
    Wait until all due foreground requests finish\;
    $b\gets$ time until next request $-\delta$\;
    \If{$i\in C$}{
        \eIf{estimated runtime of \textnormal{\textsc{TrainRest}}$(i)\le b$}{
            \textsc{TrainRest}$(C[i],y_i,p_i)$\;
            Update its runtime estimate\;
            Release $C[i]$; $i\gets i+1$\;
        }{
            Wait briefly\;
        }
        \Continue\;
    }
    $\mathcal{G}\gets\emptyset$\;
    $W\gets$ uncached batches from $i$ to $\min(M,i+L-1)$\;
    \ForEach{proxy $p$ represented in $W$}{
        $J\gets$ batch IDs in $W$ using $p$, in order\;
        $r\gets K-|C|$\;
        \lIf{$i\notin J$}{$r\gets r-1$}
        $G\gets$ the longest prefix of $J$ with size in $S$, size $\le r$, and estimated forward time $\le b$\;
        \lIf{such a $G$ exists}{$\mathcal{G}\gets\mathcal{G}\cup\{G\}$}
    }
    \If{$\mathcal{G}=\emptyset$}{
        Wait briefly; \Continue\;
    }
    \eIf{a group in $\mathcal{G}$ contains $i$}{
        Select that group as $G$\;
    }{
        Select $G\in\mathcal{G}$ with the longest estimated forward time; break ties by earliest batch ID\;
    }
    $C[G]\gets\textsc{FrozenForward}(Q[G])$\;
    Update the forward runtime estimate for $G$\;
}
\caption{Feature caching with Reorder for one epoch.}
\label{alg:caching-reorder}
\end{algorithm}


When batch $i$ is cached, its remaining training work takes priority.
Otherwise, we prepare features for it if they fit, or use the gap for a later group.
Each group uses one proxy and takes the earliest uncached batches for that proxy within the lookahead range.
We keep the largest allowed group that fits both the time budget and the cache; a group that excludes $i$ must leave one cache slot free for it.
Only feature computation can move ahead: $i$ always advances in the original order, and each cached batch is released after use.
\par
\endgroup
